\cotefile{64}
\scanpage{66}
4. Détermination de $M_{11}(4)$.

Soit $E$ une courbe elliptique sur $S$ (à car. résiduelles $\ne 2$), considérons
\[
M={}_4E \tag*{(38)}
\]
schéma en groupes sur $S$ loc isom. à $(\mathbb Z/4\mathbb Z)^2$, i.e. \deleted{syst. local} faisceau abélien \deleted{\unsure{qui|pour}} \inserted{sur} $S_{\text{ét}}$ qui est un $(\mathbb Z/4\mathbb Z)_S$ module loc. libre de rg 2.

\deleted{(39)} de sorte qu'on a
\diagram{\begin{tikzcd}[ampersand replacement=\&,column sep=small,row sep=small]
\text{(39)} \& 0 \arrow[r] \& M_0 \arrow[r] \& M \arrow[r,"\mathrm{can}"] \& M_0 \arrow[r] \& 0 \\
\& \& M \arrow[u,"\mathrm{can}"'] \arrow[ur,"2\,\mathrm{id}_M"']
\end{tikzcd}}
où
\[
M_0=M\otimes_{\mathbb Z/4\mathbb Z}\mathbb F_2\simeq {}_2E .
\]

La donnée d'un $M$ équivaut à celle d'un $M_0$, (i.e. d'un $\underline{J}$ rev. étale de degré 3 de $S$ \deleted{\unsure{($\underline J$}}
\[
\underline J\simeq M_0^\ast,\qquad
M_0=V_{\underline J}(\mathbb F_2)\simeq
\mathbb F_2^{\underline J}/\underbrace{\mathbb F_2}_{\mathrm{diag}}\big), \tag*{(40)}
\]
\uline{et} d'un torseur \inserted{$P$} sous
\[
\unsure{N|V}_{\underline J}(M_0)\simeq\underline{\mathrm{End}}_{\mathbb F_2,S}(M_0)
\simeq M_0\oplus\mathbb F_2^{\underline{\unsure{\omega|w}}} \tag*{(41)}
\]
% slanted
\margin{$\underline\omega$ est le rev. étale de degré 2 ``des orientations'' associé au rev. $\underline J$ de degré 3 de $S$}

Un épinglage de $M={}_4E$ (i.e. \deleted{d'un} la donnée d'un iso
\[
(\mathbb Z/4\mathbb Z)^2_S\simeq M\;(={}_4E) \tag*{(42)}
\]
équivaut à celle d'une \inserted{trivialisation de $\underline J$, plus d'une} section \unsure{(trivialisation)|(après translation)}
% archive sheet number: 65
